\chapter{Trwałość Danych (Wykład 8)}

W poprzednich rozdziałach, omawiając architekturę MVVM i zarządzanie stanem, operowaliśmy na danych przechowywanych w pamięci operacyjnej (RAM). Zmienne wewnątrz \texttt{ViewModel}, strumienie \texttt{StateFlow} czy kolekcje danych istnieją tak długo, jak długo żyje proces aplikacji.

W ekosystemie Androida wyróżniamy trzy główne kategorie przechowywania danych, dobierane w zależności od struktury i wolumenu informacji:

\begin{enumerate}
    \item \textbf{Pliki (File API):} Przeznaczone dla danych niestrukturalnych lub strumieni bajtów, takich jak multimedia (zdjęcia, wideo), dokumenty PDF czy pliki cache.
    \item \textbf{Relacyjne Bazy Danych (SQLite / Room):} Przeznaczone dla dużych zbiorów danych strukturalnych, wymagających złożonych zapytań, relacji i transakcyjności (np. lista zadań, historia zamówień). Temat ten zostanie omówiony szczegółowo w \textbf{Rozdziale 9}.
    \item \textbf{Magazyny Klucz-Wartość (Preferences):} Przeznaczone dla lekkich danych konfiguracyjnych, flag i prostych ustawień (np. \texttt{isDarkModeEnabled}, \texttt{sessionToken}). To na tej kategorii skupiamy się w bieżącym rozdziale.
\end{enumerate}

\section{Mechanizm SharedPreferences}

Przez wiele lat standardem przechowywania prostych danych w Androidzie była biblioteka \texttt{SharedPreferences}. Choć obecnie jest wypierana przez nowsze rozwiązania, wciąż można ją spotkać w wielu istniejących projektach.

\texttt{SharedPreferences} to API, które zarządza plikiem XML przechowywanym w prywatnym katalogu aplikacji. Dane są zapisywane w formacie pary: \textbf{Klucz} (String) $\rightarrow$ \textbf{Wartość} (Int, Boolean, String, Float, Long).

\begin{minted}[frame=single,framesep=10pt]{kotlin}
class SettingsManager(context: Context) {
    // Utworzenie/Otwarcie pliku "app_settings.xml"
    private val prefs = context.getSharedPreferences(
      "app_settings", 
      Context.MODE_PRIVATE
    )

    fun saveTheme(isDark: Boolean) {
        prefs.edit()
             .putBoolean("is_dark_mode", isDark)
             .apply() 
    }

    fun isDarkMode(): Boolean {
        return prefs.getBoolean("is_dark_mode", false)
    }
}
\end{minted}

Analizując powyższy kod, warto pochylić się nad dwoma kluczowymi elementami: rolą obiektu \texttt{Context} oraz trybem otwarcia pliku.

Metoda \texttt{getSharedPreferences()} nie jest metodą statyczną - musi zostać wywołana na instancji \texttt{Context}. Wynika to z architektury bezpieczeństwa systemu Android.

Każda aplikacja w systemie Android działa jako oddzielny użytkownik i posiada swój prywatny, chroniony katalog na dysku.
Obiekt \texttt{Context} jest \textit{oknem} na zasoby systemowe i środowisko aplikacji. To on \textit{wie}, gdzie fizycznie znajduje się katalog prywatny bieżącej aplikacji. Bez niego biblioteka nie wiedziałaby, w której ścieżce utworzyć lub odczytać plik XML.

Drugim parametrem metody jest flaga \texttt{mode}. Określa ona uprawnienia systemu plików nadawane tworzonemu plikowi XML.

\begin{itemize}
    \item \textbf{\texttt{Context.MODE\_PRIVATE} (0):} Oznacza on, że utworzony plik może być odczytany \textit{wyłącznie} przez aplikację, która go stworzyła. Gwarantuje to, że inne aplikacje zainstalowane na telefonie nie będą mogły podglądać zapisanych tam danych.
\end{itemize}

Mimo prostoty, \texttt{SharedPreferences} posiada istotne wady architektoniczne, które stały się problematyczne w nowoczesnym, reaktywnym programowaniu.

\begin{enumerate}
    \item \textbf{Synchroniczne API (Blokowanie Wątku):} Najpoważniejszą wadą jest to, że metody odczytu (np. \texttt{getString}, \texttt{getBoolean}) są wywoływane synchronicznie.

    Gdy wywołujemy \texttt{prefs.getBoolean()} na Wątku Głównym (UI Thread), aplikacja musi poczekać na odczyt danych z systemu plików. Jeśli plik preferencji jest duży lub urządzenie jest obciążone operacjami I/O, może to doprowadzić do gubienia klatek animacji (Jank), a w skrajnych przypadkach do błędu \textbf{ANR (Application Not Responding)}.

    \item \textbf{Brak natywnej reaktywności:} \texttt{SharedPreferences} nie udostępnia strumieni danych. Aby obserwować zmiany w ustawieniach, programista musi rejestrować złożony listener \texttt{OnSharedPreferenceChangeListener}. Nie integruje się to naturalnie z biblioteką Kotlin Coroutines ani Flow.
\end{enumerate}

\section{DataStore}

Google wprowadziło \textbf{DataStore} jako następcę \texttt{SharedPreferences}. Biblioteka ta została zaprojektowana od podstaw z myślą o asynchroniczności i bezpieczeństwie typów.

Główne zalety \texttt{DataStore}:
\begin{itemize}
    \item \textbf{Integracja z Korutynami i Flow:} Operacje odczytu są wystawiane jako strumień \texttt{Flow}, a zapisu jako funkcje zawieszające \texttt{suspend}.
    \item \textbf{Bezpieczeństwo wątkowe:} Operacje dyskowe są domyślnie przenoszone na \texttt{Dispatchers.IO}.
    \item \textbf{Obsługa błędów:} Posiada wbudowane mechanizmy łapania wyjątków \texttt{IOException} podczas odczytu pliku.
\end{itemize}

Istnieją dwa warianty biblioteki:
\begin{enumerate}
    \item \textbf{Preferences DataStore:} (Omówimy go teraz) Działa jak \texttt{SharedPreferences} (mapa Klucz-Wartość), nie wymaga definiowania schematu.
    \item \textbf{Proto DataStore:} Przechowuje typowane obiekty (wymaga definicji schematu przy użyciu Protocol Buffers). Zapewnia bezpieczeństwo typów, ale jest trudniejszy w konfiguracji.
\end{enumerate}

DataStore nie zwraca wartości natychmiast. Zwraca \textbf{strumień} (\texttt{Flow}), który emituje nową wartość za każdym razem, gdy dane na dysku ulegną zmianie.

Najpierw tworzymy instancję DataStore jako \textbf{Singleton} (rozszerzenie dla \texttt{Context}). Jest to kluczowe, aby uniknąć konfliktów przy jednoczesnym dostępie do pliku.

\begin{minted}[frame=single,framesep=10pt]{kotlin}
// W pliku np. DataStoreConfig.kt
val Context.dataStore: DataStore<Preferences> by preferencesDataStore(name = "user_prefs")

object UserPrefsKeys {
    // Definiujemy klucze z określeniem typu danych
    val USER_NAME = stringPreferencesKey("user_name")
    val IS_DARK_MODE = booleanPreferencesKey("is_dark_mode")
}
\end{minted}

Dobrą praktyką jest ukrycie bezpośredniego dostępu do DataStore w warstwie repozytorium.

\begin{minted}[frame=single,framesep=10pt]{kotlin}
class UserPreferencesRepository(private val context: Context) {

    // ODCZYT: Zwracamy Flow. UI będzie nasłuchiwać zmian.
    val userNameFlow: Flow<String> = context.dataStore.data
        .catch { exception ->
            // Obsługa błędów odczytu pliku
            if (exception is IOException) {
                emit(emptyPreferences())
            } else {
                throw exception
            }
        }
        .map { preferences ->
            // Pobieramy wartość lub domyślną ("")
            preferences[UserPrefsKeys.USER_NAME] ?: ""
        }

    // ZAPIS: Funkcja zawieszająca (suspend)
    suspend fun saveUserName(name: String) {
        // Metoda edit jest transakcyjna (atomowa)
        context.dataStore.edit { preferences ->
            preferences[UserPrefsKeys.USER_NAME] = name
        }
    }
}
\end{minted}

\section{Integracja z architekturą MVVM}

Posiadając warstwę repozytorium (\texttt{UserPreferencesRepository}), która wystawia dane w postaci strumienia \texttt{Flow}, możemy zintegrować je z warstwą prezentacji. Musimy przekonwertować \textit{zimny} strumień z DataStore na \textit{gorący} stan UI, dostępny natychmiast dla widoku. Wykorzystamy do tego operator \texttt{stateIn}.

Rola ViewModelu w tym procesie jest dwojaka:
\begin{enumerate}
    \item \textbf{Transformacja:} Zamienia \texttt{Flow<String>} z repozytorium na \texttt{StateFlow<String>}.
    \item \textbf{Zarządzanie cyklem życia:} Dzięki \texttt{viewModelScope}, nasłuchiwanie zmian w pliku preferencji odbywa się tylko wtedy, gdy ekran jest aktywny.
\end{enumerate}

\begin{minted}[frame=single,framesep=10pt]{kotlin}
class UserViewModel(
    private val repository: UserPreferencesRepository
) : ViewModel() {

    // 1. Stan pola tekstowego (tymczasowy, edytowalny przez użytkownika)
    var textFieldValue by mutableStateOf("")
        private set

    // 2. Stan zapisany na dysku (trwały, tylko do odczytu)
    // Operator stateIn konwertuje Flow na StateFlow
    val savedUserName: StateFlow<String> = repository.userNameFlow
        .stateIn(
            scope = viewModelScope,
            started = SharingStarted.WhileSubscribed(stopTimeoutMillis = 5000),
            initialValue = "" // Wartość początkowa zanim wczytamy plik
        )

    fun onTextFieldValueChanged(newValue: String) {
        textFieldValue = newValue
    }

    // Wywołanie zapisu (operacja asynchroniczna)
    fun onSaveClicked() {
        viewModelScope.launch {
            repository.saveUserName(textFieldValue)
            // UWAGA: Nie musimy ręcznie aktualizować `savedUserName`!
            // DataStore sam wykryje zmianę pliku i wyemituje nową wartość,
            // która automatycznie trafi do naszego StateFlow.
        }
    }
}
\end{minted}

\section{Przykład praktyczny: Ekran Ustawień}

Ostatnim elementem jest warstwa widoku (UI). Dzięki zastosowaniu \texttt{StateFlow}, nasz ekran staje się w pełni reaktywny.

\begin{minted}[frame=single,framesep=10pt]{kotlin}
@Composable
fun UserSettingsScreen(viewModel: UserViewModel = viewModel()) {
    val savedName by viewModel.savedUserName.collectAsStateWithLifecycle()

    Column(
        modifier = Modifier.fillMaxSize().padding(16.dp),
        horizontalAlignment = Alignment.CenterHorizontally,
        verticalArrangement = Arrangement.Center
    ) {
        Text(
            text = if (savedName.isBlank()) "Witaj nieznajomy!" 
                   else "Witaj, $savedName!",
        )
        
        Spacer(Modifier.height(32.dp))

        OutlinedTextField(
            value = viewModel.textFieldValue,
            onValueChange = viewModel::onTextFieldValueChanged,
            label = { Text("Wpisz swoje imię") },
            modifier = Modifier.fillMaxWidth()
        )
        
        Spacer(Modifier.height(16.dp))
        
        Button(
            onClick = { viewModel.onSaveClicked() },
            modifier = Modifier.fillMaxWidth()
        ) {
            Text("Zapisz trwale")
        }
    }
}
\end{minted}

Prześledźmy, co dzieje się po kliknięciu przycisku:
\begin{enumerate}
    \item \textbf{UI:} Użytkownik klika \textit{Zapisz}. Wywołuje się \texttt{viewModel.onSaveClicked()}.
    \item \textbf{ViewModel:} Uruchamia korutynę i zleca \texttt{repository.saveUserName()}.
    \item \textbf{DataStore:} Zapisuje dane do pliku na dysku (wątek IO).
    \item \textbf{DataStore (Callback):} Wykrywa zmianę w pliku i emituje nową wartość do strumienia \texttt{userNameFlow}.
    \item \textbf{ViewModel:} \texttt{StateFlow} odbiera nową wartość i aktualizuje swój stan.
    \item \textbf{UI:} Komponent \texttt{UserSettingsScreen} otrzymuje powiadomienie o zmianie stanu i przerysowuje się (Recomposition) z nowym imieniem.
\end{enumerate}

\section{Bezpieczeństwo wątkowe: DataStore a Dispatchers.IO}

Wróćmy na chwilę do największej bolączki \texttt{SharedPreferences} - blokowania wątku głównego. Jak pamiętamy, wywołanie \texttt{prefs.getString()} działo się dokładnie tam, gdzie zostało wywołane (często na UI Thread).

Biblioteka została zaprojektowana tak, aby można było bezpiecznie wywoływać jej metody z wątku głównego, a ciężka praca (operacje I/O, zapis na dysk) zostanie automatycznie przeniesiona na odpowiedni wątek w tle.

DataStore wewnętrznie wykorzystuje korutyny i dyspozytory. Niezależnie od tego, czy odczytujesz dane, czy je zapisujesz, biblioteka deleguje te zadania do puli wątków \texttt{Dispatchers.IO}, która jest zoptymalizowana pod kątem operacji dyskowych i sieciowych.

\begin{enumerate}
    \item \textbf{Przy zapisie (\texttt{edit}):} Funkcja jest oznaczona jako \texttt{suspend}. Gdy ją wywołujesz, DataStore zawiesza korutynę, przełącza się na wątek IO, dokonuje zapisu atomowego do pliku, a następnie wznawia korutynę na oryginalnym wątku.
    \item \textbf{Przy odczycie (\texttt{data}):} Strumień \texttt{Flow} emituje wartości, które są odczytywane z dysku w tle.
\end{enumerate}

W tym przypadku (podobnie jak w ROOM i Retrofit (kolejne rozdziały)), nie jest konieczne manualne przełączanie za pomocą \texttt{withContext}. W przypadku DataStore poniższa konstrukcja jest \textbf{zbędna}:

\begin{minted}[frame=single,framesep=10pt,fontsize=\footnotesize]{kotlin}
// Podejście Nieprawidłowe (Nadmiarowy kod)
suspend fun saveToken(token: String) {
    // DataStore i tak przełączy się na IO, więc to wywołanie jest masłem maślanym
    withContext(Dispatchers.IO) {
        context.dataStore.edit { prefs ->
            prefs[TOKEN_KEY] = token
        }
    }
}

// Podejście Prawidłowe (Main-Safe)
suspend fun saveToken(token: String) {
    // Możemy to bezpiecznie wołać z ViewModelu (Main Thread)
    context.dataStore.edit { prefs ->
        prefs[TOKEN_KEY] = token
    }
}
\end{minted}

Należy pamiętać, że o ile \textbf{odczyt z dysku} dzieje się na IO, o tyle \textbf{transformacje} danych (operatory \texttt{.map} na strumieniu Flow) wykonują się na tym wątku, na którym nastąpiła subskrypcja (kolekcjonowanie).
Dla prostych typów (\texttt{String}, \texttt{Boolean}) nie ma to znaczenia, ale przy skomplikowanym przetwarzaniu warto pamiętać o operatorze \texttt{flowOn}, omawianym w Rozdziale 6.

\section{Pełne Przykłady Implementacji}

Poniżej znajdują się kompletne kody źródłowe omawianych zagadnień.

\subsection{Przykład 1: SharedPreferences}

Ten przykład pokazuje "stary styl". Mamy tu menedżera ustawień, który operuje na pliku XML. Zauważ, że aby zaktualizować UI po zmianie przełącznika, musimy ręcznie zmienić stan w Composable (\texttt{isNotificationsEnabled = newCheckedState}), ponieważ SharedPreferences nie powiadamia nas o zmianach automatycznie.

\begin{minted}[frame=single,framesep=10pt]{kotlin}
// --- 1. SettingsManager.kt ---
class SettingsManager(context: Context) {
    // Tryb MODE_PRIVATE: tylko ta aplikacja ma dostęp do pliku
    private val prefs = context.getSharedPreferences(
      "app_settings", 
      Context.MODE_PRIVATE
    )

    companion object {
        private const val NOTIFICATIONS_KEY = "notifications_enabled"
    }

    // ZAPIS (Asynchroniczny dzięki .apply())
    fun saveNotificationsSetting(isEnabled: Boolean) {
        prefs.edit {
            putBoolean(NOTIFICATIONS_KEY, isEnabled)
            apply()
        }
    }

    // ODCZYT (Synchroniczny - uwaga na blokowanie UI przy dużych plikach!)
    fun isNotificationEnabled(): Boolean {
        return prefs.getBoolean(NOTIFICATIONS_KEY, false)
    }
}

// --- 2. Screen.kt ---
@Composable
fun LegacySettingsScreen() {
    val context = LocalContext.current
    val settingsManager = remember { SettingsManager(context) }
    
    // Stan lokalny UI
    var isNotificationsEnabled by remember { 
        mutableStateOf(settingsManager.isNotificationEnabled()) 
    }

    Column(modifier = Modifier.padding(16.dp)) {
        Text("Ustawienia (Legacy)")
        Spacer(Modifier.height(16.dp))
        
        Row(
            modifier = Modifier.fillMaxWidth(),
            horizontalArrangement = Arrangement.SpaceBetween,
            verticalAlignment = Alignment.CenterVertically
        ) {
            Text("Powiadomienia Push")
            Switch(
                checked = isNotificationsEnabled,
                onCheckedChange = { newState ->
                    // 1. Aktualizujemy UI
                    isNotificationsEnabled = newState
                    // 2. Zapisujemy w tle
                    settingsManager.saveNotificationsSetting(newState)
                }
            )
        }
    }
}
\end{minted}

\subsection{Przykład 2: DataStore (Architektura MVVM)}

\textbf{Struktura plików:}
\begin{itemize}
    \item \texttt{DataStoreConfig.kt} - Singleton instancji DataStore.
    \item \texttt{UserPreferencesRepository.kt} - Logika odczytu (Flow) i zapisu (suspend).
    \item \texttt{UserViewModel.kt} - Konwersja Flow na StateFlow.
    \item \texttt{UserScreen.kt} - Reaktywny widok.
\end{itemize}

\begin{minted}[frame=single,framesep=10pt]{kotlin}
// --- 1. DataStoreConfig.kt ---
// Delegat tworzący Singleton. Nazwa pliku to "user_prefs.preferences_pb"
val Context.dataStore by preferencesDataStore(name = "user_prefs")

// --- 2. UserPreferencesRepository.kt ---
class UserPreferencesRepository(private val context: Context) {
    
    private object Keys {
        val USER_NAME = stringPreferencesKey("user_name")
    }

    // ODCZYT: Zwracamy Flow, który emituje dane przy każdej zmianie pliku
    val userNameFlow: Flow<String> = context.dataStore.data
        .catch { exception ->
            if (exception is IOException) emit(emptyPreferences()) 
            else throw exception
        }
        .map { preferences ->
            preferences[Keys.USER_NAME] ?: "Anonim"
        }

    // ZAPIS: Funkcja zawieszająca (działa na Dispatchers.IO)
    suspend fun saveUserName(name: String) {
        context.dataStore.edit { preferences ->
            preferences[Keys.USER_NAME] = name
        }
    }
}

// --- 3. UserViewModel.kt ---
class UserViewModel(
    private val repository: UserPreferencesRepository
) : ViewModel() {

    // Stan pola tekstowego
    var inputName by mutableStateOf("")
        private set

    // Stan z dysku (to co jest zapisane)
    // stateIn zamienia zimny Flow w gorący StateFlow dostępny dla UI
    val savedName: StateFlow<String> = repository.userNameFlow
        .stateIn(
            scope = viewModelScope,
            started = SharingStarted.WhileSubscribed(5000),
            initialValue = "Wczytywanie..."
        )

    fun onInputChange(newValue: String) {
        inputName = newValue
    }

    fun saveName() {
        viewModelScope.launch {
            repository.saveUserName(inputName)
            inputName = "" // Czyścimy pole po zapisie
        }
    }
}

// --- 4. UserScreen.kt ---
@Composable
fun UserScreen(viewModel: UserViewModel = viewModel()) {
    // Obserwujemy StateFlow. Każda zmiana w pliku DataStore
    // spowoduje automatyczną rekompozycję tego widoku.
    val savedName by viewModel.savedName.collectAsStateWithLifecycle()

    Column(
        modifier = Modifier.fillMaxSize().padding(24.dp),
        horizontalAlignment = Alignment.CenterHorizontally,
        verticalArrangement = Arrangement.Center
    ) {
        Text("Witaj, $savedName!")
        
        Spacer(Modifier.height(32.dp))

        OutlinedTextField(
            value = viewModel.inputName,
            onValueChange = viewModel::onInputChange,
            label = { Text("Zmień swoje imię") },
            modifier = Modifier.fillMaxWidth()
        )
        
        Spacer(Modifier.height(16.dp))
        
        Button(
            onClick = { viewModel.saveName() },
            modifier = Modifier.fillMaxWidth()
        ) {
            Text("Zapisz w DataStore")
        }
    }
}
\end{minted}